% Part: first-order-logic
% Chapter: proof-systems
% Section: natural-deduction

\documentclass[../../../include/open-logic-section]{subfiles}

\begin{document}

\iftag{FOL}
      {\olfileid[hi]{fol}{prf}{ntd}}
      {\olfileid[hi]{pl}{prf}{ntd}}

\olsection{स्वाभाविक निगमन}

स्वाभाविक निगमन ऐसा !!{derivation}-तंत्र है जिसका उद्देश्य वास्तविक
तर्क-प्रक्रिया का प्रतिरूप देना है (विशेषतः गणितज्ञों द्वारा अपनाई जाने वाली
सुव्यवस्थित तर्क-प्रक्रिया का)। वास्तविक तर्क-प्रक्रिया कई ``स्वाभाविक''
ढाँचों से आगे बढ़ती है। उदाहरणतः, स्थिति-भेद से प्रमाण किसी वियोजनात्मक
आधारवाक्य से निष्कर्ष स्थापित करने देता है: हम दिखाते हैं कि निष्कर्ष उसके
दोनों विकल्पों में से प्रत्येक से निकलता है। अप्रत्यक्ष प्रमाण में हम यह
दिखाकर निष्कर्ष स्थापित करते हैं कि उसका निषेध विरोधाभास तक ले जाता है।
आपादन-प्रमाण में यह दिखाकर ``यदि \dots तो \dots'' रूप का दावा स्थापित किया
जाता है कि उत्तरपद, पूर्वपद से निकलता है। स्वाभाविक निगमन इन स्वाभाविक
अनुमानों में से कुछ का औपचारिक निरूपण है। हर तार्किक संयोजक और परिमाणक के साथ
दो नियम जुड़े होते हैं—एक प्रवेश-नियम और एक निरसन-नियम—और प्रत्येक नियम ऐसे
ही किसी स्वाभाविक अनुमान-ढाँचे के अनुरूप होता है। उदाहरणतः, $\Intro{\lif}$
आपादन-प्रमाण के और $\Elim{\lor}$ स्थिति-भेद से प्रमाण के अनुरूप है। विशेष रूप
से सरल नियम $\Elim{\land}$ है, जो $!A \land !B$ से~$!A$ (अथवा $!B$) का
अनुमान लगाने देता है।

स्वाभाविक निगमन को अन्य !!{derivation}-तंत्रों से अलग करने वाला एक गुण इसमें
मान्यताओं का उपयोग है। स्वाभाविक निगमन में !!^{derivation}, !!{formula}s का
वृक्ष होती है। उसके मूल पर एक !!{formula} होता है; !!{formula}s के इस वृक्ष
की ``पत्तियाँ'' ऐसे !!{formula}s से बनती हैं जिनसे निष्कर्ष व्युत्पन्न किया
जाता है। स्वाभाविक निगमन में पत्तियों पर स्थित कुछ !!{formula}s की
!!{derivation} के भीतर भूमिका होती है, किन्तु जब !!{derivation} निष्कर्ष तक
पहुँचती है, तब तक उनका उपयोग पूरा हो चुका होता है। यह वास्तविक तर्क-प्रक्रिया
की उस पद्धति के अनुरूप है जिसमें ऐसी परिकल्पनाएँ प्रस्तुत की जाती हैं जो केवल
थोड़े समय तक प्रभावी रहती हैं। उदाहरणतः, स्थिति-भेद से प्रमाण में हम दोनों
विकल्पों में से प्रत्येक को सत्य मानते हैं; आपादन-प्रमाण में पूर्वपद को सत्य
मानते हैं; और अप्रत्यक्ष प्रमाण में निष्कर्ष के निषेध को सत्य मानते हैं।
काल्पनिक मान्यताओं को प्रस्तुत करना और फिर किसी मध्यवर्ती चरण को स्थापित करने
के बाद उन्हें हटा देना स्वाभाविक निगमन की पहचान है। स्वाभाविक निगमन की
!!{derivation} की पत्तियों पर स्थित सूत्र मान्यताएँ कहलाते हैं, और कुछ
अनुमान-नियम उन्हें सक्रिय निर्भरता से ``!!{discharge}'' कर सकते हैं। उदाहरणतः,
यदि हमारे पास कोई !!{derivation} हो जो कुछ ऐसी मान्यताओं से~$!B$ देती हो
जिनमें~$!A$ शामिल है, तो $\Intro{\lif}$ नियम हमें~$!A \lif !B$ का अनुमान
लगाने और~$!A$ रूप की किसी भी मान्यता को मुक्त करने देता है। (कौन-सी मान्यता
किस अनुमान पर मुक्त होती है, इसका लेखा रखने के लिए हम अनुमान और उसके द्वारा
मुक्त की जाने वाली मान्यताओं पर एक संख्या लगाते हैं।) जो मान्यताएँ
!!{undischarged} रहती हैं, वे !!{derivation} के अंत में मिलकर निष्कर्ष की
सत्यता के लिए पर्याप्त होती हैं; इसलिए कोई !!{derivation} यह स्थापित करती है
कि उसकी !!{undischarged} मान्यताओं से उसका निष्कर्ष तार्किकतः निकलता है।

स्वाभाविक निगमन पर आधारित संबंध $\Gamma \Proves !A$ तभी और केवल तभी होता है
जब कोई !!{derivation} ऐसी हो जिसमें $!A$ वृक्ष का अंतिम !!{sentence} हो और हर
!!{undischarged} पत्ती~$\Gamma$ में हो। $!A$ स्वाभाविक निगमन में प्रमेय तभी और
केवल तभी है जब कोई !!{derivation} ऐसी हो जिसमें $!A$ अंतिम !!{sentence} हो और
सभी मान्यताएँ !!{discharged} हों। उदाहरण के लिए, यह !!{derivation} दिखाती है
कि $\Proves (!A
\land !B) \lif !A$:
\begin{prooftree}
  \AxiomC{$\Discharge{!A \land !B}{1}$}
  \RightLabel{\Elim{\land}}
  \UnaryInfC{$!A$}
  \DischargeRule{\Intro{\lif}}{1}
  \UnaryInfC{$(!A \land !B) \lif !A$}
\end{prooftree}
लेबल~$1$ बताता है कि मान्यता $!A \land !B$ को \Intro{\lif} अनुमान पर
!!{discharged} किया जाता है।
  
समुच्चय~$\Gamma$ स्वाभाविक निगमन में तभी और केवल तभी असंगत है जब
$\Gamma \Proves \lfalse$ हो। नियम \FalseInt{} के कारण असंगत समुच्चय से कोई
भी !!{sentence} !!{derive} किया जा सकता है।

स्वाभाविक-निगमन तंत्र गेरहार्ड गेन्ट्सेन और स्तानिस्वाव याश्कोव्स्की ने
1930-दशक में विकसित किए; बाद में डाग प्रावित्स और फ्रेडरिक फिच ने उन्हें
आगे बढ़ाया। इसके अनुमान प्रमाण की स्वाभाविक विधियों का प्रतिरूप हैं, इसलिए यह
दार्शनिकों में लोकप्रिय है। फिच द्वारा विकसित रूपों का उपयोग प्रायः आरंभिक
तर्कशास्त्र की पाठ्यपुस्तकों में होता है। तर्कशास्त्र के दर्शन में कभी-कभी यह
माना गया है कि स्वाभाविक निगमन के नियम तार्किक संक्रियकों के अर्थ देते हैं
(``प्रमाण-सैद्धांतिक अर्थविज्ञान'')।

\end{document}
